\section{Compressed planar capping beyond nonnegative Euler characteristic}
\label{sec:cocycle-planar}

The annulus criterion suggests checking genus rather than Euler
characteristic alone. A connected orientable surface with Euler
characteristic $-5$ can be planar with seven boundary circles. If precisely
one of those circles is essential on the exterior torus, the other six can
be capped, giving a compressing disc. The new optional native branch
recognizes this configuration using a compressed boundary census and
independent replay. It introduces no additional surface candidates.

\subsection{The elementary capping criterion}

Let $M$ be an authenticated knot exterior with boundary torus $T$, and let
$S$ be a connected orientable properly embedded surface. Write $b$ for its
number of boundary circles. Classification gives
\[
 2g(S)=2-\chi(S)-b.
\]
Suppose $\chi(S)+b=2$ and exactly one boundary circle is essential on $T$.
Every other boundary circle bounds a disc on $T$. Choose an innermost
inessential circle among those remaining. Its boundary disc has interior
disjoint from the remaining boundary curves. Attach this disc to $S$ and
push the cap slightly into a boundary collar, rounding the corner.
Properness supplies a sufficiently small collar in which this operation
avoids the rest of the surface. Repeat for all inessential circles,
choosing collar depths to preserve disjointness of the caps. This also
handles nested boundary circles; the boundary discs need not initially
have pairwise disjoint interiors on $T$.

Each cap preserves connectedness and orientability, increases Euler
characteristic by one and removes one boundary component. The resulting
surface has one boundary component and Euler characteristic
$\chi(S)+(b-1)=1$. It is therefore a disc, and its remaining boundary is
essential on $T$. The source knot is the unknot. The certificate stores
$S$ and the boundary census; it does not claim that the stored normal
surface contains a disc component or extract a normal vector for the
capped disc.

Primitivity is not an additional requirement of this topological capping
lemma. The implementation nevertheless retains its established restricted
cocycle inspector: authenticated source geometry, coherent primitive class,
and either a zero spanning tree or a checked minimum-span witness. These
supply connectedness and orientability. Euler characteristic and boundary
counts alone cannot replace those hypotheses. In particular, a disconnected
surface can satisfy a total Euler identity without being planar and
connected.

\subsection{Counting essential circles with two compressed weights}

The boundary-only normal arc system has one point for each intersection
of $\partial S$ with a boundary edge of the triangulation. Its interval
pairings follow normal arcs in boundary faces. Orbits of this system are
exactly the boundary circles. This point order differs from the full
surface edge-point order; the new geometry helper constructs its offsets
from the sorted boundary edges themselves.

Choose two cycles forming a basis of $H_1(T;\mathbb F_2)$. On the block
for a boundary edge, put a two-dimensional integer weight whose entries
indicate whether that edge belongs to each basis cycle. The weight of an
orbit counts its intersections with the two cycles. Reduction modulo two
gives its homological intersection vector. Since the intersection form of
a torus is nondegenerate, this vector vanishes exactly when the curve's
mod-two homology vanishes. An essential simple curve on a torus represents
a primitive integral class, whose reduction modulo two is nonzero.
Consequently a boundary orbit is essential exactly when at least one
of its two weight sums is odd.

The maintained weighted interval-orbit engine computes a histogram
$(w,m_w)$, where $m_w$ counts circles with integer weight vector $w$.
The two required counts are
\[
 b=\sum_w m_w,\qquad
 e=\sum_{w:\,w\bmod2\ne(0,0)}m_w.
\]
Connectedness and orientability already come from the cocycle witness,
so this query needs neither a full-surface orbit census nor its orientation
double. Only boundary orbits carry weights.
Zero-weight circles remain present in the histogram. Counts and
multiplicities stay binary. There is no loop over boundary points, normal
pieces or individual circles in discovery or replay.

For $N$ tetrahedra the boundary relation and initial weight intervals have
$O(N)$ descriptions, with weight dimension two. They use the already
accounted-for weighted-AHT transport bounds of the earlier component
sections, including polynomially many run breakpoints in the explicit
input and trace. Geometry, integer bit costs, basis construction and
independent binary rank checking remain part of the cost. This is a
polynomial encoded query on a supplied surface, not a complete search for
a genus-zero surface. The candidate family's incompleteness remains the
obstacle to a general recognition bound.

\subsection{Certificate separation and bounded optional discovery}

The new schema is \code{diagram-cocycle-planar-v1}. It contains the six
cocycle fields plus a boundary homology basis and a weighted boundary-orbit
certificate. Its independent inspector first replays the
cocycle-surface checker with its existing acceptance rules. It then builds
the boundary relation from that call's validated geometry, validates
the supplied basis by binary homology rather than repeating the producer's
tree--cotree construction, and invokes the independently implemented
weighted-orbit verifier. Only the replayed histogram supplies $b$ and $e$.
Acceptance requires $\chi+b=2$ and $e=1$.

The wrapper imports no cocycle, flow, orbit or weight-transport discovery.
The public seed verifier dispatches the new tag without changing the
acceptance rules of disc, annulus or legacy component certificates. Its
summary records genus zero and \code{unknot\_witness=planar-cap};
\code{compressing\_discs} still describes the actual stored surface.

The native search tries this query only after a candidate fails the
existing disc and enabled annulus tests. Raw, tree, optimized and optional
face candidates retain their previous order. A discovered planar proof
must pass the full independent source and boundary inspector before a
positive verdict is emitted. Negative query results supply no knottedness
claim. The first implementation repeated finite-triangulation preparation
inside the boundary producer and twice inside each positive checker.
Profiling identified this as the largest cost. The revised cocycle
inspector has a private detailed return carrying its freshly reconstructed
geometry and coordinate analysis; its public return remains the same
summary. The planar checker uses that result within the current call.
The producer likewise retains validated source geometry and its chosen
boundary basis only for the current native search. It reuses coordinate
analysis when that candidate has just been fully inspected, and otherwise
checks its new coordinate vector against the retained geometry.

No caller-supplied geometry cache is accepted by the certificate checker.
Each independent replay still reconstructs and authenticates its own
source, and checks all surface and boundary claims. This removes repeated
work without replacing positive replay by trust in discovery. All work uses the existing shared guard allowance; the seed API's
\code{max\_cycles} additionally bounds the sum of boundary orbit-discovery
cycles over all candidates. Incomplete boundary queries can fall through
to later candidates, while global work exhaustion and cancellation still
terminate the bounded stage.

The new query is explicitly requested with \code{planar=True} in the seed
API. The recognizer exposes \code{normal\_seed\_planar=True}; the CLI flag
is \code{--normal-seed-planar}. Its default is false. This preserves the
previous search records exactly when the new query is not requested.
Unlike the annulus implication, the planar test adds an orbit computation
on continuing candidates; the measurements below evaluate that cost.

\subsection{Coverage and boundary-census controls}

On the same 84 source inputs and two-million-guard allowance, the four-tree
policy increases from 24 to 25 native positives; it has 58 completed
restricted misses and one cap. With 24 tree trials, coverage increases
from 27 to 30 positives, with 53 misses and one cap. Gordian remains the
capped source. With the planar option off, all recorded results under both
policies are identical to the previous annulus audit, including work,
certificates, stages and statistics. Every previous positive is independently
replayed. Nineteen ordinary-portfolio inputs return the expected answers
under both settings without reaching the native stage.

The five-crossing \code{optimized-positive} and its repeated source
\code{random-54} now succeed at tree trial 2 on a surface with
$\chi=-5$ and seven boundary circles. Guard work falls from 54,318 to 29,111.
The four-crossing \code{random-1} succeeds at tree trial 2 with
$\chi=-4$ and six boundary circles. Its 20,856 guards replace a four-tree
miss costing 36,784, or a later annulus at trial 10 costing 44,244 under the
extended policy.

Three previously missed extended-search sources now succeed:
\code{random-25} and \code{random-176} at trial 12, and
\code{random-162} at trial 18. Their witnesses are three-boundary planar
surfaces with $\chi=-1$. The sources \code{random-25} and \code{random-176}
share a certificate; their work rises from a 43,177-guard miss to a
61,043-guard success. The source \code{random-162} rises from 75,603 to 157,897 guards.
These additional boundary queries are real work, despite adding coverage.

The seven-crossing \code{random-150} stops on its raw four-boundary
planar surface with $\chi=-2$, but its work rises from 18,281 to 25,049.
The old route reached a disc at tree trial 2 more cheaply than the raw
planar census and replay. The seven affected source labels therefore yield
five distinct planar certificates, rather than seven distinct proofs.
Each is independently replayed and checked against Regina and an expanded
boundary census.

A separate survey regenerates raw seeds, all distinct candidates from
24 tree trials, and the original span optimum for the 74 sources with at
most twelve crossings. All 1,798 supplied surfaces are connected and
orientable according to Regina. For every surface, the compressed
boundary count agrees with Regina, and the entire parity histogram agrees
with independently expanded boundary pairings. At most 230 boundary
points are expanded in these small-fixture controls. The compressed query
uses at most 118 orbit cycles on this corpus.

Exactly 345 surveyed surfaces are planar and have one essential boundary
circle. All come from known unknots and all receive independent
source-bound planar-certificate replay. Of these, 44 have negative Euler
characteristic, across fourteen source labels. Thus the extension is not
just a new tag for discs and annuli. The five witnesses actually selected
by recognition receive additional Regina and expanded-boundary checks,
for 1,803 comparisons of each kind across the survey and source audit.
These totals count comparisons, not necessarily distinct surfaces.

Focused tests cover earlier termination, producer-free replay, strict
fields, changed sources, false bases, altered weight sums and
multiplicities, positive-genus and null-class trefoil surfaces, and a
primitive disconnected trefoil control. A scaled seven-boundary fixture
has $7\cdot2^{2000}$ boundary circles and $2^{2000}$ essential circles;
the compressed census and independent verifier recover those binary
counts without expansion. Another test transports a valid source proof
through 20,000-bit tetrahedral height offsets and hexadecimal JSON.
Exact work and shared cycle boundaries, cancellation in discovery and
replay, strict Boolean options and unchanged default results are checked.

\subsection{Regression checks and the scheduling decision}

All 23 focused tests pass in 3.292 seconds. After the geometry-reuse
refactoring, all 1,210 maintained tests pass in 220.299 seconds. The final
source audit takes 75.389 seconds and the boundary survey takes 57.621
seconds. These are correctness runs, partly overlapping the full test
suite; their elapsed times are not performance comparisons. The final
audit, survey and benchmark each pin 493 source files.

The initial source and survey records, before geometry reuse, are retained.
All surveyed boundary records and all five selected certificates are
identical after refactoring. Verdicts under both source policies are
unchanged, and the off-arm records still exactly reproduce the preceding
annulus audit. The initial timing implementation is preserved as
hash-verified source snapshots, together with the earlier audit driver's
fixture map. This distinguishes the tested revisions from the final
implementation rather than assigning historical measurements its source
hashes.

Both timing runs start from source PDs and include complete recognition,
any exact fallback, independent replay of the returned native certificate,
and certificate hashing. Earlier recognition shortcuts are disabled to
reach the native stage; the seed has no local work cap, and the global
deadline remains twelve seconds. Each case has a warm-up round and five
measured rounds, randomizing off, off A/A, on and on A/A order. Tree trials
are four except for the three new extended positives, which request 24;
face discovery remains disabled. No tests, audits, builds or commits
overlap either timing run.

All 440 measured calls and 88 warm-ups across the initial and revised
implementations complete. Per-implementation verdicts and certificate
hashes agree across both runs. Native work is stable within each revision;
the intentional reductions after geometry reuse are recorded rather than
claimed unchanged. The initial run takes 28.735 seconds and the final run
22.443 seconds. The table reports the final revision only: primary-arm
median times and medians of paired ratios, which need not equal the ratio
of the displayed medians.

\begin{center}
\small\input{tables/cocycle_planar_recognize.tex}
\end{center}

The earlier seven-boundary witness improves the complete endpoint by a
paired factor of $1.121$: approximately 89 milliseconds becomes 79.
The earlier raw planar witness is near parity at 89 milliseconds.
The new four-tree positive is modestly slower, with a 0.916 off/on ratio
and noticeable off-arm A/A variation. The three new extended positives
are materially slower: their off/on ratios are 0.631, 0.530 and 0.669.
The existing optimized-annulus positive also slows because earlier
unsuccessful candidates now receive boundary queries; its ratio is 0.794.
The trefoil and genus-one continuing misses have ratios 0.912 and 0.831.
Raw disc and raw annulus controls bypass boundary discovery and show no
algorithmic speedup.

These data support an optional stronger certificate search, not a new
universal default. Geometry reuse removes much of the initial overhead,
but the extra queries and larger certificates can still cost more than
the preceding exact fallback. The maintained default therefore remains
\code{planar=False}. The whole native stage remains optional in the
ordinary portfolio, whose nineteen measured inputs never reach it.

Reproduction uses \path{../fast/normal_orbit_research/planar.py}.
Its modes are \code{audit}, \code{survey} and
\code{recognize --rounds 5}. Regina and expanded boundary enumeration
appear only in the audit controls. The incoming research inventory and
checked \code{openai/math} head \code{fd4aeeb2e} are unchanged. The new
capping proof uses elementary surface and torus topology and does not
require a zero-surgery theorem. No complete genus-zero search, explicit
normal capped-disc extraction, hierarchy cutting, or general
subexponential or quasipolynomial recognition bound is established here.



